\SourceLeaf{387}%
issue (1810) it ceased publication altogether. General \textsc{Guilleminot}, who took over as director of the Depot in 1822, first conceived the idea of continuing the journal on a regular basis. His acting deputy, General \textsc{Delachasse de Verigny}, brought the plan to maturity. It was also decided to make some changes to the format and to reprint, in the quarto format chosen for the continuation, the seven earlier issues. copies of which were out of print. The first two volumes of the new edition are intended for this reprint of the earlier issues. The first lies before us together with three volumes of the continuation. The second is to follow shortly.

The \textit{Dépôt de la guerre} was first established in 1688 under \textsc{Louis} XIV by the minister \textsc{Louvois}. Initially, however, it was only an archive where all official army correspondence was deposited. It attained its present form only gradually, as its scope widened and its organisation was consolidated. It took on its distinctive grandeur only once it became the centre that brought together all the results produced by an independent corps, the engineer-geographers. Some idea of the scale of the institution's activities can be gained from the fact that in 1801 its annual costs were estimated at 110000 francs, excluding staff salaries. In 1793 the latter amounted to 231100 francs.

High expectations are warranted for a journal that can draw on so extraordinarily rich a source. The more the editors focus on publishing important material hitherto inaccessible to the public and leave out what does nothing to advance science, the better those expectations will be met.

We find the earlier articles much further from this ideal than the more recent continuation. Indeed, apart from an essay on the hydrography of part of France and a note on the history of the \textit{Dépôt de la guerre}, the articles in the first volume are of such a kind that they could just as easily have been written had the \textit{Dépôt} never existed. Without denying the interest that several essays could have had thirty years ago and that some still have today, one can
